% chapters/chapter2/2_3_state_space.tex
% !TeX root=../../main.tex

\section{مدل فضای حالت و خطی‌سازی دینامیک سیستم}

همان‌طور که در بخش پیشین تشریح شد، در توسعه مدل تحلیلی لاگرانژ، متغیر $l(t)$ به عنوان یک سیگنال سینماتیکی تجویزشده در نظر گرفته شد و در نتیجه، معادلات حرکت برای یک مدل کاهش‌یافته با مختصات تعمیم‌یافته $x$ و $\theta$ استخراج گردید. با این وجود، به منظور استقرار الگوریتم کنترل پیش‌بین مدل غیرخطی (\lr{NMPC}) در فصل‌های آتی، معماری مدل‌سازی نیازمند بازآرایی است.

در کنترل‌کننده پیش‌بین، شتاب تغییر طول کابل ($\ddot{l}$) به عنوان یکی از متغیرهای تصمیم کنترلی انتخاب می‌شود. برای پیش‌بینی رفتار آینده سیستم و اعمال صریح قیود بر دامنه و سرعت تغییرات کابل، رابطه میان این ورودی و متغیرهای سینماتیکی کابل باید به صورت دینامیکی در مدل پیش‌بینی لحاظ شود ($\dot{l} = v_l$ و $\dot{v}_l = \ddot{l}$). از این رو، متغیرهای $l$ و $\dot{l}$ به صورت حالات افزوده (\lr{Augmented States}) به بردار وضعیت سیستم اضافه می‌شوند. این افزوده‌سازی دینامیکی حالات جهت مشتق‌گیری ورودی و تکمیل مدل پیش‌بینی انجام می‌گیرد \cite{Homaeinezhad2026Solution}. این رویکرد که در فرمول‌بندی سیستم‌های کنترل پیش‌بین برای لحاظ کردن دینامیک عملگرها و قیود ورودی استفاده می‌شود \cite{Allgower2004Nonlinear, Zak2017Introduction}، به معنای تغییر یا تناقض در معادلات لاگرانژی فصل قبل نیست؛ بلکه صرفاً تغییر فرم مدل برای تشکیل یک ساختار مداربسته و یکپارچه جهت حل مسئله بهینه‌سازی است.

\subsection{تعریف متغیرهای حالت و فرم غیرخطی فضای حالت}

بردار حالت افزوده سیستم ($X$) شامل $6$ متغیر (موقعیت و سرعت ارابه، زاویه و سرعت زاویه‌ای آونگ، و طول و سرعت کابل) و بردار ورودی مستقل ($u$) متشکل از $2$ متغیر تصمیم (نیروی محرک و شتاب کابل) به صورت زیر تعریف می‌شوند:
\begin{equation}
    X =
    \begin{bmatrix}
        x_1 \\ x_2 \\ x_3 \\ x_4 \\ x_5 \\ x_6
    \end{bmatrix}
    =
    \begin{bmatrix}
        x \\ \dot{x} \\ \theta \\ \dot{\theta} \\ l \\ \dot{l}
    \end{bmatrix},
    \qquad
    u =
    \begin{bmatrix}
        u_1 \\ u_2
    \end{bmatrix}
    =
    \begin{bmatrix}
        F \\ \ddot{l}
    \end{bmatrix}
\end{equation}
در این فرمول‌بندی، تنها دو متغیر کنترلی مستقل ($F$ و $\ddot{l}$) وجود دارند و متغیرهای $l$ و $\dot{l}$ سیگنال‌های مستقلی نیستند، بلکه از طریق دینامیک انتگرالی سیستم مستقیماً به $\ddot{l}$ وابسته‌اند ($\dot{x}_5 = x_6$ و $\dot{x}_6 = u_2$).

با استفاده از معادلات حرکت استخراج‌شده در رویکرد لاگرانژ، ماتریس اینرسی ($M$) و بردار نیروهای سمت راست معادلات ($RHS$) برحسب متغیرهای بردار حالت و ورودی جدید به شکل زیر بازنویسی می‌شوند:
\begin{equation}
    M(X) =
    \begin{bmatrix}
        M+m             & m l \cos \theta \\
        m l \cos \theta & m l^2
    \end{bmatrix}
\end{equation}
\begin{equation}
    \begin{split}
        RHS_1(X,u) ={} & F - b_c \dot{x} - m \ddot{l} \sin \theta                               \\
                       & - 2m \dot{l} \dot{\theta} \cos \theta + m l \dot{\theta}^2 \sin \theta
    \end{split}
\end{equation}
\begin{equation}
    RHS_2(X,u) = -2m l \dot{l} \dot{\theta} - m g l \sin \theta - b_p \dot{\theta}
\end{equation}

دترمینان ماتریس اینرسی به صورت تحلیلی برابر است با $|M| = m l^2(M+m\sin^2\theta)$. با توجه به اینکه $l$ اکنون یکی از متغیرهای بردار حالت ($x_5$) است، معکوس این ماتریس به فرم زیر استخراج می‌شود:
\begin{equation}
    M^{-1}(X) = \frac{1}{m l^2 (M + m \sin^2 \theta)}
    \begin{bmatrix}
        m l^2            & -m l \cos \theta \\
        -m l \cos \theta & M+m
    \end{bmatrix}
\end{equation}

با محاسبه شتاب‌های مکانیکی از رابطه $[\ddot{x}, \ddot{\theta}]^T = M^{-1}(X) [RHS_1, RHS_2]^T$ و الحاق معادلات سینماتیکی کابل، مدل پیش‌بینی فضای حالت غیرخطی سیستم ($\dot{X} = f(X,u)$) به صورت یکپارچه و دقیق استخراج می‌شود:
\begin{equation}
    \dot{X} =
    \begin{bmatrix}
        \dot{x}       \\
        \ddot{x}      \\
        \dot{\theta}  \\
        \ddot{\theta} \\
        \dot{l}       \\
        \ddot{l}
    \end{bmatrix} =
    \begin{bmatrix}
        x_2           \\
        f_x(X,u)      \\
        x_4           \\
        f_\theta(X,u) \\
        x_6           \\
        u_2
    \end{bmatrix}
\end{equation}
که در آن، توابع $f_x$ و $f_\theta$ توصیف‌گر شتاب‌های غیرخطی سیستم هستند. این دستگاه معادلات $6$ حالته، مبنای تحلیلی الگوریتم بهینه‌ساز پیش‌بین در فصل‌های آتی است. جزئیات محاسبات نمادین (\lr{Symbolic}) و راستی‌آزمایی این معادلات جهت تکرارپذیری، در کد \ref{code:state_space} از پیوست \ref{app:matlab_codes} ارائه شده است.

\subsection{خطی‌سازی و استخراج ماتریس‌های ژاکوبین}

برای تحلیل رفتار محلی سیستم و بررسی ساختار درونی آن، نقطه کار (\lr{Equilibrium Point}) متناظر با وضعیت سکون کامل سیستم در نظر گرفته می‌شود. در این وضعیت، ارابه در مبدأ متوقف است ($x=0$)، آونگ به صورت قائم آویزان است ($\theta=0$)، طول کابل در مقدار ثابت $l_0$ قرار دارد و تمامی سرعت‌ها و نیروهای محرک برابر صفر هستند:
\begin{equation}
    X_{eq} = \begin{bmatrix} 0 \\ 0 \\ 0 \\ 0 \\ l_0 \\ 0 \end{bmatrix}, \quad
    u_{eq} = \begin{bmatrix} 0 \\ 0 \end{bmatrix}
\end{equation}

با اعمال بسط تیلور حول نقطه کار و محاسبه تحلیلی ماتریس‌های ژاکوبین نسبت به بردار حالت افزوده ($A = \frac{\partial f}{\partial X}\big|_{eq}$) و بردار ورودی ($B = \frac{\partial f}{\partial u}\big|_{eq}$)، ماتریس‌های سیستم خطی‌شده با ابعاد $6 \times 6$ و $6 \times 2$ به دست می‌آیند:
\begin{equation}
    \renewcommand{\arraystretch}{1.5}
    A =
    \begin{bmatrix}
        0 & 1                 & 0                     & 0                           & 0 & 0 \\
        0 & -\frac{b_c}{M}    & \frac{mg}{M}          & \frac{b_p}{M l_0}           & 0 & 0 \\
        0 & 0                 & 0                     & 1                           & 0 & 0 \\
        0 & \frac{b_c}{M l_0} & -\frac{g(M+m)}{M l_0} & -\frac{b_p(M+m)}{M m l_0^2} & 0 & 0 \\
        0 & 0                 & 0                     & 0                           & 0 & 1 \\
        0 & 0                 & 0                     & 0                           & 0 & 0
    \end{bmatrix}
\end{equation}

\begin{equation}
    \renewcommand{\arraystretch}{1.5}
    B =
    \begin{bmatrix}
        0                & 0 \\
        \frac{1}{M}      & 0 \\
        0                & 0 \\
        -\frac{1}{M l_0} & 0 \\
        0                & 0 \\
        0                & 1
    \end{bmatrix}
\end{equation}

\subsection{تحلیل کنترل‌پذیری و کوپلینگ مرتبه اول}

بررسی ساختار ماتریس‌های خطی‌شده $A$ و $B$، رفتار دینامیکی تعاملات زیرسیستم‌ها را در مجاورت وضعیت تعادل فاش می‌کند. همان‌طور که از این ماتریس‌ها استنتاج می‌شود، سیستم در تقریب خطی به دو زیرسیستم مجزا تفکیک می‌یابد:
\begin{enumerate}
    \item \textbf{زیرسیستم مکانیکی:} شامل متغیرهای $[x, \dot{x}, \theta, \dot{\theta}]^T$ که در این همسایگی محلی تنها از طریق ورودی $u_1 = F$ تحریک می‌شود.
    \item \textbf{زیرسیستم کابل:} شامل متغیرهای افزوده $[l, \dot{l}]^T$ که یک زنجیره انتگرالی مرتبه دوم تشکیل داده و مستقیماً توسط ورودی $u_2 = \ddot{l}$ هدایت می‌شود.
\end{enumerate}

در ماتریس توزیع ورودی ($B$)، درایه‌های متناظر با اثرگذاری ورودی شتاب کابل ($\ddot{l}$) بر شتاب ارابه و آونگ مطلقاً صفر هستند. این پدیده به معنای عدم وجود کوپلینگ خطی مرتبه اول (\lr{First-Order Decoupling}) میان ورودی تغییر طول کابل و دینامیک نوسان در نقطه کار قائم است. با این حال، این ویژگی استقلال، کاملاً محلی بوده و صرفاً به مدل خطی‌شده پیرامون وضعیت قائم تعلق دارد. در مدل غیرخطی اصلی، کوپلینگ‌های قدرتمندی نظیر شتاب کوریولیس ($-2m \dot{l} \dot{\theta} \cos \theta$) و مؤلفه اینرسی مکانیزم وینچ ($-m \ddot{l} \sin \theta$) حضور دارند که وابستگی مستقیم و لاینفک نوسان بار به حرکت کابل را در زوایای غیرصفر تضمین می‌کنند.

هدف اساسی از توسعه مدل افزوده $6$ حالته در این پژوهش، دقیقاً حفظ همین کوپلینگ‌های غیرخطی در بستر معماری کنترل‌کننده است. با تزریق این مدل جامع به الگوریتم، کنترل‌کننده توانایی آن را می‌یابد که علاوه بر مدیریت دو ورودی کنترلی محدود ($F$ و $\ddot{l}$)، محدودیت‌های فیزیکی مستقیم بر متغیرهای حالت کابل، نظیر کران جابجایی ($l_{\min} \le l \le l_{\max}$)، را به عنوان بخشی تفکیک‌ناپذیر از مسئله بهینه‌سازی در افق پیش‌بینی ارضا نماید.